% !TeX root = ../main.tex
\section{Discussion and Future Work}
\label{sec:discussion}
\paragraph{What connects the modules.}
The main coordination mechanism is explicit in the targets: P learns the
outcome of G's current choice, and I learns the benefit of the selected P/G
continuation. This aligns preparation with the policy that actually follows.
The matched inspection comparisons show the clearest practical benefit:
selective observation maintains task outcomes while reducing travel and time.
The continuation studies also reveal a branch tradeoff. For example, one
continuation checkpoint reaches $26/36$ direct validation correctness but reduces the
inspected branch from $35/36$ to $30/36$. Evaluating both branches during
selection protects the useful inspected continuation.

\paragraph{Current operating range.}
The experiments use nine familiar templates, one mobile training seed,
ideal base servos, known-geometry navigation, and a calibrated handle
estimator with color, axis, and radius priors. Four stations and two grasps
define the mobile candidate library.

\paragraph{Next experiments.}
Three comparisons follow directly from the results. First, hold the candidate
library and controller fixed while comparing current-G-conditioned P labels
with independently defined station labels. Second, replace estimated handle
geometry and depth exclusions separately with geometric reference inputs to
measure how much error each interface contributes. Third, expand one action
dimension at a time and add recovery triggered by low joint margin or contact
loss. Multiple training seeds and frozen new-template tests would then measure
how these changes carry beyond the current scene families.
